%%%%%%%%%%%%%%%%%%%%%%%%%%%%%%%%%%%%%%%%%%%%%%%%%%%%%%%%%%%
% Definition aller wichtigen Variablen und Einstellungen! %
%%%%%%%%%%%%%%%%%%%%%%%%%%%%%%%%%%%%%%%%%%%%%%%%%%%%%%%%%%%

% Angaben Student Allgemein

% Autor
\newcommand{\baauthor}{Sebastian Meissner}
% Matrikelnummer: Ohne führende 0 und ES, ohne Versionsnummer 
\newcommand{\bamatnum}{3084104}
% Titel der Bachelorarbeit
\newcommand{\batitle}{Anomalieerkennung in adaptiven Systemen mit Generativer KI}
% Fachsemester
\newcommand{\basem}{14}
% Ort (für Titelseite und eidesstattliche Erklärung)
\newcommand{\balocation}{Dorsten}

% Adresse Titelseite
\newcommand{\baadress}{
\baauthor \\
Wachtelstraße 74 \\
46282
\balocation \\
Matrikelnummer: \bamatnum
}

% Angaben Lehrstuhl (Betreuer, Gutachter, Zweitgutachter)

% Betreuer
\newcommand{\babetreuer}{Prof.\,Dr.\,Andreas Metzger}
% Erstgutachter
\newcommand{\bagutachter}{Prof.\,Dr.\,Andreas Metzger}
% Zweitgutachter
\newcommand{\bazgutachter}{Prof.\,Dr.\,ABC}


% Meta-Daten für das erzeugte PDF-File, Keywords bei Bedarf hinzufügen
\hypersetup{
	pdftitle    = {Seminararbeit \baauthor},
	pdfauthor   = {\baauthor},
	pdfsubject  = {\batitle},
	%pdfkeywords = {Stichwort1, Stichwort2, Stichwort3,....}
}

% Settings

% Abstand über Kapitelname verringern
\renewcommand*{\chapterheadstartvskip}{\vspace*{.5\baselineskip}}

% Numerierungstiefe
\setcounter{secnumdepth}{2}
\setcounter{tocdepth}{2}

% Pagestyle ,Header und Fußleisten

% Linien unter der Kopf- und über der Fußzeile
\KOMAoptions{
   headsepline = true,
   footsepline = true
}


% Vorherige Headereinstellungen überschreiben
\pagestyle{scrheadings}
\clearscrheadfoot

% Entfernt Nummerierung von Kapitelname in Kopfzeile
\renewcommand*{\chaptermarkformat}{}

% Aktueller Kapitelname wird ausgelesen und in der Kopfzeile links innen angezeigt
\automark[chapter]{chapter}
\ihead[\headmark]{\headmark}

% Links innen BA-Thema, Rechts außen BA-Thema Seitenzahl (1mm Abstand an beiden Seiten)
\ofoot*{\pagemark\hspace*{0.1mm}}

%%%%%%%%%%%%%%%%%%%%%%%%%%%%%%%%%%%%%%%%%%%%%%%%%%%%%%%%%%
% Wichtig! Sollte der Titel besonders lang sein und über %
% zwei Zeilen gehen: Variable \batitle ersetzen mit dem  %
% Themennamen und an geeigneter Stelle mit \\ einen      %
% manuellen Zeilenumbruch erzwingen.					 %
%%%%%%%%%%%%%%%%%%%%%%%%%%%%%%%%%%%%%%%%%%%%%%%%%%%%%%%%%%
\ifoot*{\hspace*{0.1mm}\footnotesize{\batitle}}

% Schriftgöße in Kopfleiste
\setkomafont{pageheadfoot}{\sffamily\small}
\renewcommand{\headfont}{\small}

% Schriftart zu Times New Roman
\setkomafont{disposition}{\bfseries}

% Schriftgröße für Chapter, Section, etc. ändern
\addtokomafont{chapter}{\LARGE}
\addtokomafont{section}{\Large}
\addtokomafont{subsection}{\large}
\addtokomafont{paragraph}{\itshape}

% Einrückung links bei jedem neuen Kapitel, Section, Subsection, Subsubsection
\setlength{\parindent}{0.5cm}

% Nummerierungsart der Bildunterschriften von Abbildungen und Tabellen
\makeatletter 
\counterwithout{figure}{chapter}
\counterwithout{table}{chapter}

\renewcommand{\thefigure}{\arabic{figure}} 
\renewcommand{\thetable}{\arabic{table}} 
\renewcommand{\cfttabpresnum}{\bfseries Tabelle } 
\renewcommand{\cftfigpresnum}{\bfseries Abbildung } 
\renewcommand{\cftfigaftersnum}{:} 
\renewcommand{\cfttabaftersnum}{:} 
\setlength{\cftfignumwidth}{2,5cm}                     
\setlength{\cfttabnumwidth}{2,0cm}                                           
\setlength{\cftfigindent}{0cm}                                                     
\setlength{\cfttabindent}{0cm} 
\makeatother


%%%%%%%%%%%%%%%%%%%%%%
% Nützliche Commands %
%%%%%%%%%%%%%%%%%%%%%%

% Leere Seite einfügen
\newcommand*{\blankpage}{%
  \clearpage\thispagestyle{empty}\null\newpage%
}

% Neuen Absatz einfügen mit 0.5cm Abstand nach links
% Nutzung:  ....abc def. \nAbsatz ghi jkl....
% Ergebnis: 
%			....abc def.
%			
%				ghi jkl....
\newcommand{\nAbsatz}{\newline\newline \noindent\hspace*{0.5cm}}

% Neuen Absatz beginnen am Anfang der neuen Zeile Abstand von 0.5cm nach links einfügen
\newcommand{\nEinrueck}{\noindent\hspace*{0.5cm}}

% z.B. mit Leerzeichen (kleiner und großer erster Buchstabe)
\newcommand{\stzB}{z.\,B. }
\newcommand{\stZB}{Z.\,B. }

% d.h. mit Leerzeichen (kleiner und großer erster Buchstabe)
\newcommand{\stdh}{d.\,h. }
\newcommand{\stDh}{D.\,h. }

% u.a. mit Leerzeichen (kleiner und großer erster Buchstabe)
\newcommand{\stua}{u.\,a. }
\newcommand{\stUa}{U.\,a. }

% Bild einfügen mit Bildunterschrift, Label für interne Verlinkung und Skalierung
% Nutzung: \mImage{caption}{label}{width-factor}{path}
% BSP: \mImage{Abbildung ABC}{abb:abc}{0.5}{uni-due-logo.png}
%
% Anmerkung: Die maximale Breite (bei Scalefaktor 1.0) ist die Breite des restlichen
% Textes. Bei mehrzeiliger Bildunterschrift wird der Text zentriert dargestellt.
\newcommand{\mImage}[4]{
	\begin{figure}[h!]
		\centering
		\includegraphics[width=\textwidth * \real{#3}]{./Bilder/#4}
		\captionsetup{width=\textwidth * \real{#3}}
		\caption{\label{#2}#1}
	\end{figure}    
}
\newcommand{\mImageCfg}[5]{
	\begin{figure}[#5]
		\centering
		\includegraphics[width=\textwidth * \real{#3}]{./Bilder/#4}
		\captionsetup{width=\textwidth * \real{#3}}
		\caption{\label{#2}#1}
	\end{figure}    
}

% Referenz auf Abbildung, referenzierbar anhand Label des zuvor definierten Bildcommands
% Zeigt: ".. Abbildung 2.x.x zeigt.."
% Nutzung: ".. \mImageRef{abb:uni-logo} zeigt.."
\newcommand{\mImageRef}[1]{\hyperref[#1]{Abbildung~\ref*{#1}}}

% Referenz auf Seitenzahl einer Abbildung, referenzierbar anhand Label des zuvor definierten Bildcommands
% Zeigt: "Auf Seite 12 sieht man..."
% Nutzung: "Auf \mPageRef{abb:uni-logo} sieht man..."
\newcommand{\mPageRef}[1]{\hyperref[#1]{Seite~\pageref*{#1}}}